\documentclass[12pt]{extarticle}

% ===== Encodage & langue =====
\usepackage{fontspec}
\usepackage{polyglossia}
\setdefaultlanguage{french}

% ===== Police Verdana (compiler avec XeLaTeX ou LuaLaTeX) =====
\setsansfont{Verdana}
\renewcommand{\familydefault}{\sfdefault}

% ===== Interligne =====
\usepackage{setspace}
\setstretch{1.1}

% ===== Marges réduites =====
\usepackage{geometry}
\geometry{a4paper, top=1.8cm, bottom=1.8cm, left=2cm, right=2cm}

% ===== Mathématiques =====
\usepackage{amsmath}

% ===== Espacement réduit autour des équations =====
\setlength{\abovedisplayskip}{4pt}
\setlength{\belowdisplayskip}{4pt}
\setlength{\abovedisplayshortskip}{2pt}
\setlength{\belowdisplayshortskip}{2pt}

% ===== Espacement réduit des sections =====
\usepackage{titlesec}
\titlespacing*{\section}{0pt}{1.2ex}{0.6ex}

% ===== Titre du document =====
\title{\textbf{Correction du contrôle séquence 1 version 5}}
\author{}
\date{}

\begin{document}
\maketitle

% --------------------------------------------------
\section*{Exercice 1}
\begin{align*}
16 - 7 - 1 + 24 - 4 &= 9 - 1 + 24 - 4 \\
                     &= 8 + 24 - 4 \\
                     &= 32 - 4 \\
                     &= 28
\end{align*}

% --------------------------------------------------
\section*{Exercice 2}

\begin{align*}
\textbf{a)}
9 + 4 \times 5 &= 9 + 20 \\
               &= 29
\end{align*}

\textbf{b)}
\begin{align*}
47 - 9 \times 4 &= 47 - 36 \\
                &= 11
\end{align*}

\textbf{c)}
\begin{align*}
35 + 14 \div 7 &= 35 + 2 \\
               &= 37
\end{align*}

\textbf{d)}
\begin{align*}
9 + 7 \times 4 + 6 - 8 \div 2 &= 9 + 28 + 6 - 4 \\
                              &= 39
\end{align*}

% --------------------------------------------------
\section*{Exercice 3}
\begin{align*}
5 \times 6 + 2 \times 9 &= 30 + 18 \\
                        &= 48
\end{align*}

La dépense est de 48~€.

% --------------------------------------------------
\section*{Exercice 4}

\textbf{a)} $34 - (8 + 6) \div 2$
\begin{align*}
&= 34 - 14 \div 2 \\
&= 34 - 7 \\
&= 27
\end{align*}

\textbf{b)} $(7 + 6) \times 5 - 9 + 4$
\begin{align*}
&= 13 \times 5 - 9 + 4 \\
&= 65 - 9 + 4 \\
&= 60
\end{align*}

\textbf{c)} $(27 - (7 + 9)) \times 4$
\begin{align*}
&= (27 - 16) \times 4 \\
&= 11 \times 4 \\
&= 44
\end{align*}

\textbf{d)} $6 \times (26 - (9 + 11)) - 7$
\begin{align*}
&= 6 \times (26 - 20) - 7 \\
&= 6 \times 6 - 7 \\
&= 36 - 7 \\
&= 29
\end{align*}

% --------------------------------------------------
\section*{Exercice 5}
\begin{align*}
17 - 5 \times 3 &= 17 - 15 \\
                &= 2
\end{align*}

Une gomme coûte 2~€.

% --------------------------------------------------
\section*{Exercice bonus 1}

\textbf{a)} $7 \times 5 + 15$

\textbf{b)} $6 \times (9 + 4)$

% --------------------------------------------------
\section*{Exercice bonus 2}

\textbf{a)} La somme du produit de 7 par 6 et de 5.

\textbf{b)} Le produit de la somme de 5 et de 8 par 6.

\end{document}